\section{Addition and multiplication of cycles}\label{sec:operations}

The coproduct $\Delta$ in $H$ is defined by pullback under cycle
addition; the operation $\rho_d$ is pullback under $D\mapsto dD$,
for $d\geq1$ \cite[Definitions~3.9 and~3.12]{Pardon}. The two
arguments of the coproduct are cycles on independent families of
threefolds. In the Borel construction this distinction introduces
the class of the diagonal in a product of two projective spaces.

\begin{proposition}\label{prop:operations}
For a fixed balanced pair $(C,L)$, write $a_m=a_m(C,L)$. Then
\begin{equation}\label{eq:operations}
 \Delta a_m=\sum_{s+t=m}a_s\otimes a_t,\qquad
 \rho_d a_m=
 \begin{cases}a_{m/d}&d\mid m,\\0&d\nmid m.\end{cases}
\end{equation}
\end{proposition}
\begin{proof}
The preimage of $Z_m$ under multiplication by $d$ is empty unless
$d\mid m$. Otherwise pull the tuple of Lemma~\ref{lem:support}
back to $Z_{m/d}$. The pulled-back functions have the same weights,
with common zero $(m/d)[C]$. Naturality of the relative Thom class
and Lemma~\ref{lem:stabilization} give the formula for $\rho_d$.

For the coproduct fix $s+t=m$. Over the diagonal in $B_N^2$, addition
gives a map from the common-base family of $Z_s\times Z_t$ to
$Z_{m,N}$. In the two-base description of $\Delta$
\cite[(4.17)--(4.21)]{Pardon}, the $(s,t)$ component is represented
by pulling back $u_N(F)$ along addition and then applying the
oriented Gysin map $\delta_{N!}$ over this diagonal. Its target is
the independent product $Z_{s,N}\times Z_{t,N}$. We can take
$|F|=N$ with common zero $m[C]$. The tuple pulled back along
addition has common zero $(s[C],t[C])$, since both summands are
effective cycles.

Choose compactly supported cut-downs given by tuples $F_s,F_t$ on
the two factors, of lengths $n_s,n_t$, and set $M=n_s+n_t$. On
the common-base family their concatenation also has compact common
zero set. The proof of Lemma~\ref{lem:stabilization}, applied to
the diagonal action on $Z_s\times Z_t$, identifies its class at
stage $M$ with the pulled-back class at stage $N$ after
stabilization. We may make this substitution before applying the
diagonal Gysin map. Indeed, for the simultaneous hyperplane
inclusion $i$ the identity
\[
 (i\times i)_!\,\delta_{N!}=\delta_{N+1,!}\,i_!
\]
follows by composition of oriented embeddings on the bases.
Hyperplane pushforward in each of the two independent bases
preserves the image of the corresponding relative class.

At stage $M$ the concatenated class on the diagonal is the
restriction of $u_M(F_s)\boxtimes u_M(F_t)$. By the projection
formula its diagonal pushforward is this external product times
\[
 [\delta_M(B_M)]=\sum_{i+j=M}h_1^ih_2^j .
\]
The two factors of the $(i,j)$ summand have relative cohomological
degrees $2(n_s+i-M)$ and $2(n_t+j-M)$, whose sum is zero.
A positive degree here is a negative virtual dimension. Its image
in the semi-Fano algebra vanishes by \cite[Theorem~1.2]{Pardon},
as in the proof of \cite[Theorem~7.6]{Pardon}. Consequently only
$i=n_t$, $j=n_s$ survives. Multiplication by these powers of the
hyperplane classes completes each tuple to length $M$ at stage
$M$, giving $a_s\otimes a_t$ with coefficient one.
An empty factor admits the empty tuple and the same argument, with
any extra copies of $h$ represented by the Euler class of
$\cO(1)$. Summation over $s+t=m$ proves \eqref{eq:operations}.
\end{proof}
